\section{Introduction}
\label{sec:introduction}

Decentralized finance has generated a large public archive of incidents, exploit writeups, protocol post-mortems, and security dashboards, yet those artifacts rarely line up into a single evidentiary chain. Corpus-wide datasets are broad but noisy; protocol audits are deep but local; exploit proofs of concept are concrete but narrow. As a result, DeFi security discourse often shifts between incompatible units of analysis without explicitly stating when the underlying evidence has changed. A repository may contain hundreds of scanner findings, a dashboard may report thousands of incidents, and a post-mortem may present a single elegant exploit path, but none of those views alone answers a central research question: which publicly reported incidents correspond to replayable pure on-chain exploit mechanisms rather than operational compromise, social engineering, or context-specific failures?

This paper addresses that gap using the artifact set that already exists in the repository. The repository contains a normalized incident corpus derived from the SlowMist archive, five deep audits of high-value protocols, a 24-protocol study of pure on-chain attack surfaces, and a growing replay suite for historically important incidents. Instead of treating these assets as separate workstreams, we assemble them into a staged workflow inspired by research automation systems such as ARA. The purpose is not to claim that every incident has already been reproduced. The purpose is to demonstrate how broad incident evidence can be progressively filtered, constrained, and promoted into mechanism-level claims with explicit confidence boundaries.

Our empirical thesis is deliberately conservative. Public DeFi incident volume overstates replayable exploitability for mature protocols. The reason is not that incident feeds are inaccurate, but that they aggregate distinct phenomena: direct contract exploits, flash-loan-assisted logic failures, oracle manipulation, governance abuse, bridge-message abuse, private-key compromise, wallet theft, scams, and operational leakage. Once those classes are normalized and compared against protocol-level verification, the set of high-confidence pure on-chain attack paths becomes much smaller. Within the current repository artifacts, oracle manipulation, flash-loan-assisted logic abuse, and callback-driven reentrancy remain the most replayable classes, whereas many other reported incidents either depend on off-chain compromise or remain only conditionally exploitable.

The paper makes three contributions. First, it presents an explicit incident-to-replay workflow that links corpus construction, protocol screening, replay engineering, and paper synthesis. Second, it provides a quantitative summary of the repository evidence: 2{,}036 incidents, 1{,}378 DeFi-labeled cases, 87 manually triaged audit findings across five high-value protocols, and a 24-protocol pure on-chain verification study. Third, it distills these materials into a set of claims that are stronger than a scanner report yet more disciplined than an overconfident exploit catalog. The result is not a definitive census of DeFi exploitability; it is a reproducible foundation for a more methodologically coherent research program.
